% Iteration 8 change manifest as two structured cards: chg-1 hardens the
% publish-state shell guard, chg-2 makes execution-risk warnings salient on the
% next model turn and adds contract-aware heuristics. Both ship together at the
% iteration-7 boundary as a deliberate keep-and-improve on the architecture
% installed by iterations 5 and 6.
\begin{figure}[t]
    \centering
    \scriptsize
    \begin{tcbraster}[raster columns=2,
                      raster equal height,
                      raster column skip=4pt,
                      raster row skip=0pt,
                      raster left skip=0pt,
                      raster right skip=0pt,
                      boxrule=0.6pt,
                      left=4pt,right=4pt,top=3pt,bottom=3pt,
                      coltitle=white]
        \begin{tcolorbox}[colback=aheSky!18,colframe=aheNavy,colbacktitle=aheNavy,
                          title={\scriptsize\textbf{chg-1，迭代 8，commit \texttt{ca35f53}，层级：工具实现}}]
            \textbf{\textcolor{aheNavy}{涉及文件}} \\
            \rule{\linewidth}{0.3pt}\\[1pt]
            $\bullet$ \texttt{tools/shell\_tools/run\_shell\_command.py} \\[3pt]
            \rule{\linewidth}{0.3pt}\\
            \textbf{\textcolor{aheNavy}{改动内容}} \\
            \rule{\linewidth}{0.3pt}\\[1pt]
            把两类软性理由升级为硬阻断。删除任何非 \texttt{/tmp} 的受保护输出，现在是硬阻断。把任何非 \texttt{/tmp} 的受保护根目录重置为空状态，同样是硬阻断。\texttt{ALLOW\_POST\_SUCCESS\_RESET} 令牌对其他类别的成功后互锁依然存在，但它已不能再抹掉已验证的实际交付物，也不能清空实际的根目录。 \\[3pt]
            \rule{\linewidth}{0.3pt}\\
            \textbf{\textcolor{aheNavy}{修复的失败模式}} \\
            \rule{\linewidth}{0.3pt}\\[1pt]
            在部署检查成功之后，智能体附上覆写令牌，删除 \texttt{/git/www/hello.html} 并重置 \texttt{/git/server/refs/heads/master}，理由是``回到一个干净的仓库''；随后验证器拿到 404。 \\[3pt]
            \rule{\linewidth}{0.3pt}\\
            \textbf{\textcolor{aheNavy}{预测修复的任务}} \\
            \rule{\linewidth}{0.3pt}\\[1pt]
            2 个任务。例如：\texttt{configure-git-webserver}、\texttt{git-multibranch}。
        \end{tcolorbox}
        \begin{tcolorbox}[colback=aheTeal!12,colframe=aheBlue,colbacktitle=aheBlue,
                          title={\scriptsize\textbf{chg-2，迭代 8，commit \texttt{a4a4a29}，层级：中间件}}]
            \textbf{\textcolor{aheBlue}{涉及文件}} \\
            \rule{\linewidth}{0.3pt}\\[1pt]
            $\bullet$ \texttt{workspace/middleware/execution\_risk\_hints.py} \\[3pt]
            \rule{\linewidth}{0.3pt}\\
            \textbf{\textcolor{aheBlue}{改动内容}} \\
            \rule{\linewidth}{0.3pt}\\[1pt]
            新增了一个 \texttt{BeforeModelHook}，把上一步产生的任何执行风险提示提升为 FRAMEWORK 提醒，显示在下一个模型轮次的最上方，使警告进入推理上下文，而不是尾随在工具输出之后。新增了按任务一次性的契约推断，用于识别整洁布局或单文件交付类契约，以及官方封装或指定修订版本类契约。新增了两条工具调用之后的启发式规则：当智能体在整洁布局的实际目录树内编译或构建时发出警告；当契约指明使用官方封装、而命令却改用裸的 \texttt{SentenceTransformer} 或 \texttt{AutoModel} 风格 API 时发出警告。 \\[3pt]
            \rule{\linewidth}{0.3pt}\\
            \textbf{\textcolor{aheBlue}{修复的失败模式}} \\
            \rule{\linewidth}{0.3pt}\\[1pt]
            迭代 6 的中间件发出了正确的警告，但只写进了工具输出；智能体往往在下一个模型轮次才做出发布或停止的决定，并忽略了这些警告。显著性提升加上契约感知的启发式规则弥合了这一缺口。 \\[3pt]
            \rule{\linewidth}{0.3pt}\\
            \textbf{\textcolor{aheBlue}{预测修复的任务}} \\
            \rule{\linewidth}{0.3pt}\\[1pt]
            4 个任务。例如：\texttt{polyglot-c-py}、\texttt{polyglot-rust-c}、\texttt{mteb-retrieve}、\texttt{pytorch-model-recovery}。
        \end{tcolorbox}
    \end{tcbraster}
    \caption{在迭代 7 边界处一并写下、并作为迭代 8 的 harness 上线的两条变更清单条目。\texttt{chg-1} 强化了已有的发布态 shell 守卫，使覆写令牌不再能抹掉已验证的实际交付物。\texttt{chg-2} 让执行风险警告在下一个模型轮次无法被忽视，并补充了两条契约感知的启发式规则。二者的作用范围都是刻意划定的：\texttt{chg-1} 阻止破坏性命令本身，\texttt{chg-2} 修补迭代 6 中间件的显著性缺口。}
    \label{fig:manifest-hardblock}
\end{figure}
